\documentclass[11pt,a4paper]{article}
\usepackage[T1]{fontenc}
\usepackage{lmodern}
\usepackage[margin=26mm]{geometry}
\usepackage{amsmath,amssymb,amsthm,booktabs,array,longtable}
\usepackage{microtype,multicol}
\usepackage[hidelinks]{hyperref}
\setlength{\parskip}{3pt}
\setlength{\emergencystretch}{2em}
\raggedbottom
\allowdisplaybreaks[1]
\clubpenalty=10000
\widowpenalty=10000
\newtheorem{theorem}{Theorem}[section]
\newtheorem{lemma}[theorem]{Lemma}
\newtheorem{corollary}[theorem]{Corollary}
\theoremstyle{remark}
\newtheorem{remark}[theorem]{Remark}
\newcommand{\Z}{\mathbb Z}
\newcommand{\Q}{\mathbb Q}
\newcommand{\N}{\mathbb N}
\newcommand{\F}{\mathbb F}
\newcommand{\ps}[1]{#1[[x]]}
\newcommand{\Rev}{\operatorname{Rev}}
\newcommand{\rev}{^{\langle-1\rangle}}
\newcommand{\Plus}{\operatorname{Plus}}
\newcommand{\Minus}{\operatorname{Minus}}
\newcommand{\coeff}[2]{[x^{#1}]#2}
\newcommand{\E}{\mathcal E}
\Urlmuskip=0mu plus 1mu
\hypersetup{pdftitle={A380056: Exponential coefficients and congruences modulo four and five},pdfauthor={Dev Goyal},pdfsubject={Separate manuscript with AI disclosure and Lean proof map}}
\title{A380056\\[4pt]\large Exponential coefficients and congruences modulo four and five}
\author{Dev Goyal}
\date{24 October 2026}
\begin{document}
\maketitle
\begin{center}
\fbox{\begin{minipage}{0.9\linewidth}\small
\textbf{AI disclosure.} ChatGPT (OpenAI) generated the mathematical arguments,
Lean proofs, and this manuscript in response to Dev Goyal's requests.
Dev Goyal is the sole human author. The source sequence and conjecture or
observation are credited to Paul D. Hanna.
\textbf{Independent human review is pending.}
\end{minipage}}
\end{center}
\begin{abstract}
If $\sum a_nx^n/n!=(e^x-1)/\cos(2x)$, then $a_n\equiv1\pmod4$ for $n\ge1$ and $a_n\equiv2\cdot2^n+3\cdot4^n\pmod5$ for all $n\ge0$. Thus $5\mid a_n$ exactly when $4\mid n$. Integer binomial convolution constructs the sequence and proves its literal rational exponential-generating-function identity.
\end{abstract}

\section{Source, notation, and scope}

The source is Paul D. Hanna's OEIS entry A380056 \cite{oeis056},
retrieved on 24 October 2026. Its initial sequence index is 1.
We use the auxiliary value $a_0=0$ and distinguish the exponential
coefficients $a_n=n![x^n]A$ from ordinary rational coefficients.

All series are formal. The set $\N$ includes zero, and $[x^n]U$ denotes a
coefficient. A normalized series is $x+O(x^2)$; $O(x^r)$ means divisibility
by $x^r$, over any coefficient ring used below. We write $U\circ V=U(V(x))$
only when $V(0)=0$. Compositional inversion, denoted $\Rev(U)$ or $U\rev$,
is distinct from multiplicative inversion of a unit. Congruences are
coefficientwise; ordinary products and compositional iterates are distinct.

\section{A380056: exponential coefficients and binomial convolution}

Here the source is the exponential generating function
\cite{oeis056}
\[
 A(x)=\sum_{n\ge0}a_n\frac{x^n}{n!}
     =\frac{e^x-1}{\cos(2x)},\qquad a_0=0. \tag{1}
\]
The integer sequence must be distinguished from the ordinary rational
coefficients of this formal series. We first establish that distinction
and the source equation exactly.

For two sequences over a commutative ring, define binomial convolution by
\[
 (f\star g)_n=\sum_{k=0}^n\binom nk f_k g_{n-k}.
\]
Write $\delta=(1,0,0,\ldots)$ and $d=(0,1,1,\ldots)$, and define
\[
 c_{2j}=(-4)^j,\qquad c_{2j+1}=0.
\]

\begin{lemma}[Integral source construction]\label{lem:egfsource}
There is a unique integer sequence $a$ satisfying $a\star c=d$. Its
exponential generating function is exactly (1).
\end{lemma}

\begin{proof}
Since $c_0=1$, set $a_0=0$ and, for $n>0$, set
\[
 a_n=1-\sum_{k=0}^{n-1}\binom nk a_k c_{n-k}.
\]
This determines integers without division and gives the stated
convolution identity. Induction proves uniqueness over any commutative
coefficient ring.

Over $\Q$, let $\E(f)=\sum_{n\ge0}f_nx^n/n!$. The identity
$\binom nk k!(n-k)!=n!$ gives
$\E(f\star g)=\E(f)\E(g)$ coefficientwise. Also
$\E(c)=\cos(2x)$ and $\E(d)=e^x-1$. Thus
$\E(a)\cos(2x)=e^x-1$. The cosine has constant coefficient one and is
a multiplicative unit, so its inverse gives (1), uniquely in $\Q[[x]]$.
By construction $n![x^n]\E(a)=a_n$.
\end{proof}

\begin{theorem}[The two A380056 conjectures]\label{thm:egf}
For every $n\ge1$, $a_n\equiv1\pmod4$. For every $n\ge0$,
\[
 a_n\equiv2\cdot2^n+3\cdot4^n\pmod5. \tag{2}
\]
Consequently the residues modulo five at index classes $0,1,2,3$ modulo
four are respectively $0,1,1,3$, and
\[
 5\mid a_n\quad\Longleftrightarrow\quad4\mid n.
\]
In particular $5\mid a_{4n}$ for every positive $n$.
\end{theorem}

\begin{proof}
All positive-index coefficients of $c$ vanish modulo four, so $c=\delta$
in binomial convolution over $\Z/4\Z$. Reduction of $a\star c=d$
therefore gives $a=d$, proving the first assertion.

For the second, work with sequences over $\F_5$ and put $E_s(n)=s^n$,
with $0^0=1$. The binomial theorem gives
\[
 E_s\star E_t=E_{s+t}.
\]
These are divided-power coefficient sequences. No factorial denominators
or ordinary exponential series over $\F_5$ are needed. Since $-4=1$ in
$\F_5$, the sequence $c$ is one at even indices and zero at odd indices,
and hence
\[
 c=3(E_1+E_4),\qquad d=E_1-E_0.
\]
For $b=2E_2+3E_4$, bilinearity gives
\[
 \begin{aligned}
 b\star c
 &=6E_3+6E_1+9E_0+9E_3\\
 &=E_1-E_0=d
 \quad\text{over }\F_5.
 \end{aligned}
\]
The recurrence in Lemma~\ref{lem:egfsource} remains uniquely solvable
over $\F_5$, so $a=b$ after reduction, proving (2). Both $2^4$ and $4^4$
equal one modulo five. The four residues follow by evaluating the four
index classes, and only class zero gives zero.
\end{proof}

The endpoints in \texttt{Round4.ExpCosine} prove the integer convolution
recurrence, its uniqueness, the factorial coefficient interpretation,
the literal quotient in (1), uniqueness among rational formal series,
both conjectures, and the stronger complete classification modulo five.
The source bridge uses Mathlib's formal cosine series rescaled by two.


\section{Separate Lean proof and checked scope}

The entry file is \nolinkurl{lean/Entries/A380056.lean}. Its explicit
theorem statements and proof terms use the namespace
\nolinkurl{FinalCollection.A380056}. The complete imported implementations
and their shared mathematical dependencies are included in the collection.
The following local theorem names give the proof map:

\begin{multicols}{2}\raggedright\begin{itemize}\setlength{\itemsep}{2pt}\setlength{\parsep}{0pt}
\item \nolinkurl{integer_egf_definition}.
\item \nolinkurl{source_egf}.
\item \nolinkurl{literal_source_equation}.
\item \nolinkurl{literal_source_unique}.
\item \nolinkurl{integer_coefficients}.
\item \nolinkurl{first_source_conjecture}.
\item \nolinkurl{second_source_conjecture}.
\item \nolinkurl{exact_mod_five_pattern}.
\item \nolinkurl{exact_divisibility_by_five}.
\item \nolinkurl{mod_five_exponential_formula}.
\item \nolinkurl{mod_five_period}.
\end{itemize}\end{multicols}

The main implementation module is \nolinkurl{Round4.ExpCosine}.
The catalogue records the full original endpoint names and the new entry
theorems. The project pins Lean 4.19.0 and Mathlib revision
\nolinkurl{c44e0c8ee63ca166450922a373c7409c5d26b00b}. The collection
report records source recompilation and transitive axiom checks. Accepted
axioms are propositional extensionality, classical choice, and quotient
soundness. All index conclusions above are unbounded; exact finite
coefficient checks in the package provide separate corroboration.

\paragraph{Limits and status.}
These are exponential coefficients $n![x^n]A$, not ordinary rational coefficients. The literal rational EGF and its uniqueness are formally proved. Modular convolution does not divide by factorials in finite characteristic.
The manuscript consolidates completed results; it does not assert global
novelty, exclusive priority, peer review, OEIS acceptance, or a new submission.
Standard algebraic ingredients are not claimed as new discoveries.

\begin{thebibliography}{9}
\raggedright
\bibitem{oeis056}
Paul D. Hanna, \emph{OEIS A380056}, 28 January 2025.
\url{https://oeis.org/A380056}. Retrieved 24 October 2026.
\end{thebibliography}
\end{document}
